\documentclass[12pt]{article}
\usepackage[utf8]{inputenc}
\usepackage[T1]{fontenc}
\usepackage{hyperref}
\title{Mandara Language Resources for K-12 Curriculum Development}
\author{Compiled by Claude Sonnet 5 (claude-sonnet-5) for Andr\'as Kornai}
\date{}
\begin{document}
\maketitle

\section{Introduction}
Mandara (ISO 639-3 tbf; also called Tabar) is an Austronesian language of the Tabar Group in New Ireland Province, Papua New Guinea. About 4,000 speakers are cited. Three dialects are listed: Simberi, Tatau and Tabar. It is unrelated to the Chadic Mandara language of Africa. Tok Pisin and English are used for wider communication.

\section{Primary Educational Resources}
No school materials were found. A Wikipedia summary mentions a written form and a Bible, but the claim is undated and I did not verify it.

\section{Linguistic and Reference Materials}
\begin{itemize}
\item Wikipedia article on the Mandara language. \url{https://en.wikipedia.org/wiki/Mandara_language}
\item Wikipedia article on the Tabar Group. \url{https://en.wikipedia.org/wiki/Tabar_Group}
\end{itemize}

\section{Religious and Community Resources}
\begin{itemize}
\item Joshua Project page for the Tabar people. It is general missions-oriented information. \url{https://joshuaproject.net/people_groups/15156/PP}
\end{itemize}

\section{Digital and Audio Resources}
\begin{itemize}
\item Global Recordings Network page for Mandara (tbf). The search result shows only the page title, so I did not confirm which recordings it lists. \url{https://globalrecordings.net/en/language/tbf}
\end{itemize}

\section{Curriculum Development Context}
Little documentation turned up. A curriculum would need to establish the orthography and any existing scripture translation with local church and community bodies, and would need new materials.
\end{document}